\section{Quasi-finite maps}
\label{section-quasi-finite}

\noindent
Consider a ring map $R \to S$ of finite type.
A map $\Spec(S) \to \Spec(R)$ is quasi-finite
at a point if that point is isolated in its fibre.
This means that the fibre is zero dimensional at that point.
In this section we study the basic properties of this
important but technical notion. More advanced material
can be found in the next section.

\begin{lemma}
\label{lemma-isolated-point}
Let $k$ be a field.
Let $S$ be a finite type $k$-algebra.
Let $\mathfrak q$ be a prime of $S$.
The following are equivalent:
\begin{enumerate}
\item $\mathfrak q$ is an isolated point of $\Spec(S)$,
\item $S_{\mathfrak q}$ is finite over $k$,
\item there exists a $g \in S$, $g \not\in \mathfrak q$ such that
$D(g) = \{ \mathfrak q \}$,
\item $\dim_{\mathfrak q} \Spec(S) = 0$,
\item $\mathfrak q$ is a closed point of $\Spec(S)$ and
$\dim(S_{\mathfrak q}) = 0$, and
\item the field extension $\kappa(\mathfrak q)/k$ is finite
and $\dim(S_{\mathfrak q}) = 0$.
\end{enumerate}
In this case $S = S_{\mathfrak q} \times S'$ for some
finite type $k$-algebra $S'$. Also, the element $g$
as in (3) has the property $S_{\mathfrak q} = S_g$.
\end{lemma}

\begin{proof}
Suppose $\mathfrak q$ is an isolated point of $\Spec(S)$, i.e.,
$\{\mathfrak q\}$ is open in $\Spec(S)$.
Because $\Spec(S)$ is a Jacobson space (see
Lemmas \ref{lemma-finite-type-field-Jacobson} and
\ref{lemma-jacobson})
we see that $\mathfrak q$ is a closed point. Hence
$\{\mathfrak q\}$ is open and closed in $\Spec(S)$.
By
Lemmas \ref{lemma-disjoint-decomposition} and
\ref{lemma-disjoint-implies-product} we may
write $S = S_1 \times S_2$ with $\mathfrak q$
corresponding to the only point $\Spec(S_1)$.
Hence $S_1 = S_{\mathfrak q}$ is a zero dimensional
ring of finite type over $k$. Hence it is finite over $k$
for example by Lemma \ref{lemma-Noether-normalization}.
We have proved (1) implies (2).

\medskip\noindent
Suppose $S_{\mathfrak q}$ is finite over $k$.
Then $S_{\mathfrak q}$ is Artinian local, see
Lemma \ref{lemma-finite-dimensional-algebra}. So
$\Spec(S_{\mathfrak q}) = \{\mathfrak qS_{\mathfrak q}\}$ by
Lemma \ref{lemma-artinian-finite-length}.
Consider the exact sequence $0 \to K \to S \to S_{\mathfrak q}
\to Q \to 0$. It is clear that $K_{\mathfrak q} = Q_{\mathfrak q} = 0$.
Also, $K$ is a finite $S$-module as $S$ is Noetherian and
$Q$ is a finite $S$-module since $S_{\mathfrak q}$ is finite over $k$.
Hence there exists $g \in S$, $g \not \in \mathfrak q$ such that
$K_g = Q_g = 0$. Thus $S_{\mathfrak q} = S_g$ and
$D(g) = \{ \mathfrak q \}$. We have proved that (2) implies (3).

\medskip\noindent
Suppose $D(g) =  \{ \mathfrak q \}$. Since $D(g)$ is open by
construction of the topology on $\Spec(S)$ we see that
$\mathfrak q$ is an isolated point of $\Spec(S)$.
We have proved that (3) implies (1).
In other words (1), (2) and (3) are equivalent.

\medskip\noindent
Assume $\dim_{\mathfrak q} \Spec(S) = 0$. This means that
there is some open neighbourhood of $\mathfrak q$ in $\Spec(S)$
which has dimension zero. Then there is an open neighbourhood of the
form $D(g)$ which has dimension zero. Since $S_g$ is Noetherian
we conclude that $S_g$ is Artinian and
$D(g) = \Spec(S_g)$ is a finite discrete set, see
Proposition \ref{proposition-dimension-zero-ring}.
Thus $\mathfrak q$ is an isolated point of $D(g)$ and,
by the equivalence of (1) and (2) above applied to
$\mathfrak qS_g \subset S_g$, we see that
$S_{\mathfrak q} = (S_g)_{\mathfrak qS_g}$ is finite over $k$.
Hence (4) implies (2). It is clear that (1) implies (4).
Thus (1) -- (4) are all equivalent.

\medskip\noindent
Lemma \ref{lemma-dimension-closed-point-finite-type-field}
gives the implication (5) $\Rightarrow$ (4).
The implication (4) $\Rightarrow$ (6) follows from
Lemma \ref{lemma-dimension-at-a-point-finite-type-field}.
The implication (6) $\Rightarrow$ (5) follows from
Lemma \ref{lemma-finite-residue-extension-closed}.
At this point we know (1) -- (6) are equivalent.

\medskip\noindent
The two statements at the end of the lemma we saw during the
course of the proof of the equivalence of (1), (2) and (3) above.
\end{proof}

\begin{lemma}
\label{lemma-isolated-point-fibre}
\begin{slogan}
Equivalent conditions for isolated points in fibres
\end{slogan}
Let $R \to S$ be a ring map of finite type.
Let $\mathfrak q \subset S$ be a prime lying over
$\mathfrak p \subset R$. Let $F = \Spec(S \otimes_R \kappa(\mathfrak p))$
be the fibre of $\Spec(S) \to \Spec(R)$, see
Remark \ref{remark-fundamental-diagram}.
Denote $\overline{\mathfrak q} \in F$ the point corresponding to
$\mathfrak q$. The following are equivalent
\begin{enumerate}
\item $\overline{\mathfrak q}$ is an isolated point of $F$,
\item $S_{\mathfrak q}/\mathfrak pS_{\mathfrak q}$ is finite over
$\kappa(\mathfrak p)$,
\item there exists a $g \in S$, $g \not \in \mathfrak q$ such that
the only prime of $D(g)$ mapping to $\mathfrak p$ is $\mathfrak q$,
\item $\dim_{\overline{\mathfrak q}}(F) = 0$,
\item $\overline{\mathfrak q}$ is a closed point of $F$ and
$\dim(S_{\mathfrak q}/\mathfrak pS_{\mathfrak q}) = 0$, and
\item the field extension $\kappa(\mathfrak q)/\kappa(\mathfrak p)$
is finite and $\dim(S_{\mathfrak q}/\mathfrak pS_{\mathfrak q}) = 0$.
\end{enumerate}
\end{lemma}

\begin{proof}
Note that $S_{\mathfrak q}/\mathfrak pS_{\mathfrak q} =
(S \otimes_R \kappa(\mathfrak p))_{\overline{\mathfrak q}}$.
Moreover $S \otimes_R \kappa(\mathfrak p)$ is of finite type over
$\kappa(\mathfrak p)$.
The conditions correspond exactly to the conditions of
Lemma \ref{lemma-isolated-point}
for the $\kappa(\mathfrak p)$-algebra $S \otimes_R \kappa(\mathfrak p)$
and the prime $\overline{\mathfrak q}$, hence they are equivalent.
\end{proof}

\begin{definition}
\label{definition-quasi-finite}
Let $R \to S$ be a finite type ring map.
Let $\mathfrak q \subset S$ be a prime.
\begin{enumerate}
\item If the equivalent conditions of Lemma \ref{lemma-isolated-point-fibre}
are satisfied then we say $R \to S$ is {\it quasi-finite at $\mathfrak q$}.
\item We say a ring map $A \to B$ is {\it quasi-finite}
if it is of finite type and quasi-finite at all primes of $B$.
\end{enumerate}
\end{definition}

\begin{lemma}
\label{lemma-quasi-finite-above-p}
Let $R \to S$ be a finite type ring map
and $\mathfrak p$ be a prime ideal of $R$.
Then the following are equivalent:
\begin{enumerate}
\item $R \to S$ is quasi-finite at all primes of $S$ lying over $\mathfrak p$,
\item $S \otimes_R \kappa(\mathfrak p)$ is a finite
$\kappa(\mathfrak p)$-algebra, and
\item $\Spec(S \otimes_R \kappa(\mathfrak p))$ is a finite set.
\end{enumerate}
\end{lemma}

\begin{proof}
Condition (1) says the topology on
$\Spec(S \otimes_R \kappa(\mathfrak p))$ is discrete (as every point
is open). Hence the equivalence of (1), (2), (3) follows from
Lemma \ref{lemma-finite-type-algebra-finite-nr-primes}.
\end{proof}

\begin{lemma}
\label{lemma-quasi-finite}
Let $R \to S$ be a finite type ring map. Then $R \to S$ is quasi-finite
if and only if for all primes $\mathfrak p \subset R$ the ring
$S \otimes_R \kappa(\mathfrak p)$ is finite over $\kappa(\mathfrak p)$.
\end{lemma}

\begin{proof}
Follows immediately from the more general
Lemma \ref{lemma-quasi-finite-above-p}.
\end{proof}

\begin{lemma}
\label{lemma-quasi-finite-local}
Let $R \to S$ be a finite type ring map. Let $\mathfrak q \subset S$
be a prime lying over $\mathfrak p \subset R$. Let
$f \in R$, $f \not \in \mathfrak p$ and $g \in S$, $g \not \in \mathfrak q$.
Then $R \to S$ is quasi-finite at $\mathfrak q$ if and only if
$R_f \to S_{fg}$ is quasi-finite at $\mathfrak qS_{fg}$.
\end{lemma}

\begin{proof}
The fibre of $\Spec(S_{fg}) \to \Spec(R_f)$ is homeomorphic
to an open subset of the fibre of $\Spec(S) \to \Spec(R)$.
Hence the lemma follows from part (1) of the equivalent conditions of
Lemma \ref{lemma-isolated-point-fibre}.
\end{proof}

\begin{lemma}
\label{lemma-four-rings}
Let
$$
\xymatrix{
S \ar[r] & S' & &
\mathfrak q \ar@{-}[r] & \mathfrak q' \\
R \ar[u] \ar[r] &  R' \ar[u] & &
\mathfrak p \ar@{-}[r] \ar@{-}[u] & \mathfrak p' \ar@{-}[u]
}
$$
be a commutative diagram of rings with primes as indicated.
Assume $R \to S$ of finite type, and $S \otimes_R R' \to S'$ surjective.
If $R \to S$ is quasi-finite at $\mathfrak q$, then
$R' \to S'$ is quasi-finite at $\mathfrak q'$.
\end{lemma}

\begin{proof}
Write $S \otimes_R \kappa(\mathfrak p) = S_1 \times S_2$
with $S_1$ finite over $\kappa(\mathfrak p)$ and such that
$\mathfrak q$ corresponds to a point of $S_1$ as in
Lemma \ref{lemma-isolated-point}. This product decomposition
induces a corresponding product decomposition for any
$S \otimes_R \kappa(\mathfrak p)$-algebra. In particular,
we obtain $S' \otimes_{R'} \kappa(\mathfrak p') = S'_1 \times S'_2$.
Because $S \otimes_R R' \to S'$ is surjective the canonical map
$(S \otimes_R \kappa(\mathfrak p)) \otimes_{\kappa(\mathfrak p)}
\kappa(\mathfrak p') \to S' \otimes_{R'} \kappa(\mathfrak p')$
is surjective and hence
$S_i \otimes_{\kappa(\mathfrak p)} \kappa(\mathfrak p') \to S'_i$
is surjective. It follows that $S'_1$ is finite over $\kappa(\mathfrak p')$.
The map $S' \otimes_{R'} \kappa(\mathfrak p') \to
\kappa(\mathfrak q')$ factors through $S_1'$
(i.e.\ it annihilates the factor $S_2'$)
because the map $S \otimes_R \kappa(\mathfrak p) \to
\kappa(\mathfrak q)$ factors through $S_1$
(i.e.\ it annihilates the factor $S_2$). Thus
$\mathfrak q'$ corresponds to a point of
$\Spec(S_1')$ in the disjoint union decomposition
of the fibre: $\Spec(S' \otimes_{R'} \kappa(\mathfrak p'))
= \Spec(S_1') \amalg \Spec(S_2')$, see
Lemma \ref{lemma-spec-product}.
Since $S_1'$ is finite over a field, it is Artinian ring,
and hence $\Spec(S_1')$ is a finite discrete set.
(See Proposition \ref{proposition-dimension-zero-ring}.)
We conclude $\mathfrak q'$ is isolated in its fibre as
desired.
\end{proof}

\begin{lemma}
\label{lemma-quasi-finite-composition}
A composition of quasi-finite ring maps is quasi-finite.
\end{lemma}

\begin{proof}
Suppose $A \to B$ and $B \to C$ are quasi-finite ring maps. By
Lemma \ref{lemma-compose-finite-type}
we see that $A \to C$ is of finite type.
Let $\mathfrak r \subset C$  be a prime of $C$ lying over
$\mathfrak q \subset B$ and $\mathfrak p \subset A$. Since $A \to B$ and
$B \to C$ are quasi-finite at $\mathfrak q$ and $\mathfrak r$ respectively,
then there exist $b \in B$ and $c \in C$ such that $\mathfrak q$ is
the only prime of $D(b)$ which maps to $\mathfrak p$ and similarly
$\mathfrak r$ is  the only prime of $D(c)$ which maps to $\mathfrak q$.
If $c' \in C$ is the image of $b \in B$, then $\mathfrak r$ is the only
prime of $D(cc')$ which maps to $\mathfrak p$.
Therefore $A \to C$ is quasi-finite at $\mathfrak r$.
\end{proof}

\begin{lemma}
\label{lemma-quasi-finite-base-change}
Let $R \to S$ be a ring map of finite type.
Let $R \to R'$ be any ring map. Set $S' = R' \otimes_R S$.
\begin{enumerate}
\item The set
$\{\mathfrak q' \mid R' \to S' \text{ quasi-finite at }\mathfrak q'\}$
is the inverse image of the corresponding set of $\Spec(S)$
under the canonical map $\Spec(S') \to \Spec(S)$.
\item If $\Spec(R') \to \Spec(R)$ is surjective,
then $R \to S$ is quasi-finite if and only if $R' \to S'$ is quasi-finite.
\item Any base change of a quasi-finite ring map is quasi-finite.
\end{enumerate}
\end{lemma}

\begin{proof}
Let $\mathfrak p' \subset R'$ be a prime lying over $\mathfrak p \subset R$.
Then the fibre ring $S' \otimes_{R'} \kappa(\mathfrak p')$ is the
base change of the fibre ring $S \otimes_R \kappa(\mathfrak p)$
by the field extension $\kappa(\mathfrak p) \to \kappa(\mathfrak p')$.
Hence the first assertion follows from the invariance of dimension
under field extension
(Lemma \ref{lemma-dimension-at-a-point-preserved-field-extension})
and Lemma \ref{lemma-isolated-point}.
The stability of quasi-finite maps under base change follows from
this and the stability of finite type property under base change.
The second assertion follows
since the assumption implies that given a prime $\mathfrak q \subset S$ we can
find a prime $\mathfrak q' \subset S'$ lying over it.
\end{proof}

\begin{lemma}
\label{lemma-quasi-finite-permanence}
Let $A \to B$ and $B \to C$ be ring homomorphisms such that $A \to C$
is of finite type. Let $\mathfrak r$ be a prime of $C$ lying over
$\mathfrak q \subset B$ and $\mathfrak p \subset A$.
If $A \to C$ is quasi-finite at $\mathfrak r$, then
$B \to C$ is quasi-finite at $\mathfrak r$.
\end{lemma}

\begin{proof}
Observe that $B \to C$ is of finite type
(Lemma \ref{lemma-compose-finite-type})
so that the statement makes sense.
Let us use characterization (3) of Lemma \ref{lemma-isolated-point-fibre}.
If $A \to C$ is quasi-finite at $\mathfrak r$, then
there exists some $c \in C$ such that
$$
\{\mathfrak r' \subset C \text{ lying over }\mathfrak p\} \cap D(c) =
\{\mathfrak{r}\}.
$$
Since the primes $\mathfrak r' \subset C$ lying over $\mathfrak q$
form a subset of the primes $\mathfrak r' \subset C$ lying over
$\mathfrak p$ we conclude $B \to C$ is quasi-finite at $\mathfrak r$.
\end{proof}

\begin{lemma}
\label{lemma-generically-finite}
Let $R \to S$ be a ring map of finite type.
Let $\mathfrak p \subset R$ be a minimal prime.
Assume that there are at most finitely many primes of $S$
lying over $\mathfrak p$. Then there exists a
$g \in R$, $g \not \in \mathfrak p$ such that the
ring map $R_g \to S_g$ is finite.
\end{lemma}

\begin{proof}
Let $x_1, \ldots, x_n$ be generators of $S$ over $R$.
By Lemma \ref{lemma-quasi-finite-above-p} the assumption means that
$S \otimes_R \kappa(\mathfrak p) = S_{\mathfrak p}/\mathfrak pS_{\mathfrak p}$
is a finite $\kappa(\mathfrak p)$-algebra.
Thus we may find monic polynomials
$P_i \in R_{\mathfrak p}[X]$ such that $P_i(x_i)$ maps to zero
in $S_{\mathfrak p}/\mathfrak pS_{\mathfrak p}$.
Since $\mathfrak p$ is a minimal prime,
$\mathfrak pR_{\mathfrak p}$ is a locally nilpotent ideal, see
Lemma \ref{lemma-minimal-prime-reduced-ring}.
Hence $\mathfrak pS_{\mathfrak p}$ is a locally nilpotent ideal, see
Lemma \ref{lemma-locally-nilpotent}. Thus
there exist $e_i \geq 1$ such that $P(x_i)^{e_i} = 0$
in $S_{\mathfrak p}$. Let $g_1 \in R$, $g_1 \not \in \mathfrak p$
be an element such that $P_i$ has coefficients in $R[1/g_1]$ for all $i$.
Next, let $g_2 \in R$, $g_2 \not \in \mathfrak p$ be an element
such that $P(x_i)^{e_i} = 0$ in $S_{g_1g_2}$. Setting $g = g_1g_2$
we win.
\end{proof}







